%lesson 4, 12/03/2026

\subsubsection{Spettri energetici limitati inferiormente}
Il caso di spettri energetici non limitati inferiormente non ha significato fisico, in quanto il sistema
tenderà sempre a raggiungere lo stato fondamentale (quando possibile): se questo ha energia "infinitamente
negativa", cioè se $Sp(H)$ non è limitato inferiormente, la particella "cade in un buco" e, per osservarla,
sarebbe necessaria energia infinita.

Si cerca ora di determinare se esistano potenziali non limitati inferiormente che però generino spettri
limitati: un esempio ben noto è il potenziale di Coulomb $V(r) = \frac{-e}{4\pi\epsilon_0}\frac{1}{r}$, che
genera lo spettro energetico dell'atomo di idrogeno, limitato inferiormente dall'energia di Rydberg.

Si consideri il caso di una particella in 3 dimensioni soggetta a un potenziale radiale del tipo
$$V(r) = \frac{-\alpha}{r^n}$$
e si supponga che esista uno stato legato per $r\to0$ e che la particella si trovi in un intervallo
dell'ordine di $r_0\sim\Delta r$: per il principio di indeterminazione si ha
$\Delta p \sim p \sim \frac{\hbar}{r_0}$ e dunque
$$\langle H \rangle = \frac{1}{2m}\langle p^2\rangle - \frac{\alpha}{r_0^n} \sim 
\frac{\hbar^2}{2mr_0^2} - \frac{\alpha}{r_0^n}$$
Intuitivamente, se $\langle H \rangle$ ha un minimo per qualche $r_0$, lo spettro è limitato inferiormente:
$$\pdv[]{}{r_0}\langle H \rangle = \frac{-\hbar^2}{mr_0^3} - \frac{n\alpha}{r_0^{n+1}} = 0 \to
r_0^{n-2} = \frac{-mn\alpha}{\hbar^2}
$$
che non ha soluzione per $n=2$; la derivata seconda $\partial^2_{r_0} \langle H \rangle$ è positiva in 
$r_0 = \frac{-mn\alpha}{\hbar^2}$ se e solo se $n=1$.

Dunque solo i potenziali con $n=1$ hanno significato fisico.

Un'intuizione per comprendere il perché esistano spettri limitati con potenziali non limitati è la
seguente: quando la particella rischia di "cadere nel buco" e la sua distanza dall'origine diminuisce, essa
si agita sempre di più per il principio di Heisenberg e, se l'attrazione del potenziale non è troppo forte,
l'energia cinetica e potenziale raggiungono un equilibrio; sono proprio l'esistenza di $\hbar$ e
dell'indeterminazione che permettono questo equilibrio, cioè la stabilità della materia.

\subsection{Lo stato fondamentale}
Prese hamiltoniane con spettro limitato inferiormente, lo stato fondamentale (o vuoto), cioè quello con
energia $E_0 = \min \{Sp(H)\}$, ha alcune importanti proprietà.

Si introducono prima due lemmi, utili per le dimostrazioni successive

\noindent\fbox{%
    \parbox{\textwidth}{%
	\textbf{Lemma 1} Se il potenziale $V(x)$ è reale, allora le autofunzioni di $H$ possono essere
	prese reali.

    \textbf{Dimostrazione} $H\psi(x) = E\psi(x) \to H\psi^*(x)=E\psi^*(x)$, le combinazioni
    \begin{align*}
	\psi_+(x) = \frac{1}{2}(\psi(x)+\psi^*(x)) = \Re(\psi) \\
	\psi_-(x) = \frac{-i}{2}(\psi(x)-\psi^*(x)) = \Im(\psi)
    \end{align*}
    sono autofunzioni reali di $H$. \textbf{QED}
    }%
}


\noindent\fbox{%
    \parbox{\textwidth}{%
    \textbf{Lemma 2} La funzione d'onda dello stato fondamentale $\psi_0(x)$ di energia $E_0$
    non ha mai zeri.

    \textbf{Dimostrazione (1D)} per assurdo, suppongo che $\exists x_0$ tale che $\psi_0(x_0)=0$: 
    questo zero deve essere semplice, poichè in caso contrario $\psi_0(x)\equiv0$ e la soluzione non è fisica;
    
    dunque $\eval{\dv[]{\psi_0}{x}}_{x_0}\neq 0$ e vale che 
    $$\begin{cases}
	\psi_0(x)>0 &\text{se }x<x_0\\
    \psi_0(x)<0 &\text{se }x>x_0
    \end{cases}
    $$
    o viceversa.

    Si consideri ora $\varphi(x):=|\psi_0(x)|$:

    $$\bra{\varphi}H\ket{\varphi} = \sumint dE dE'\bra{\varphi}\ket{E}\bra{E}H\ket{E'}\bra{E'}\ket{\varphi}
    = \sumint dE E |\bra{\varphi}\ket{E}| > E_0 \sumint dE |\bra{\varphi}\ket{E}| = E_0 $$
    dove si è usata la completezza dello spettro di $H$ nel primo e nell'ultimo passaggio.
    
    Si ha inoltre
    $$\bra{\varphi}H\ket{\varphi} = \frac{1}{2m} ||p\ket{\varphi}||^2 + \bra{\varphi}V\ket{\varphi} = 
       \int_{}^{}dx \frac{\hbar^2}{2m}|\varphi'(x)|^2 + V(x)|\varphi(x)|^2
    $$
ma, per definizione di $\varphi(x)$, $|\varphi(x)| = |\psi_0(x)|$ e $|\varphi'(x)| = |\psi_0'(x)|$, quindi
$$\bra{\varphi}H\ket{\varphi} = E_0$$
Queste affermazioni sono in contraddizione, quindi $\nexists x_0$ \textbf{QED}.
}%
}

Si enuncia ora un importante teorema:

\noindent\fbox{%
    \parbox{\textwidth}{%
    \textbf{Teorema} Lo stato fondamentale non è mai degenere

    \textbf{Dimostrazione} per assurdo, suppongo che esistano $\psi_1,\psi_2$ autostati linearmente
  indipendenti di $H$ con autovalore $E_0$: fissato $x_0\in\mathbb{R}$, considero la combinazione lineare

  $$ \psi(x) =N [\psi_1(x_0)\psi_2(x) - \psi_2(x_0)\psi_1(x)]$$
    che ha sempre energia $E_0$; o il punto $x_0$ è uno zero di $\psi(x)$, che è assurdo per il Lemma 2,
    o la funzione è identicamente nulla, che significa che $\psi_1\propto\psi_2$, che è assurdo perché 
    $\psi_1,\psi_2$ sono linearmente indipendenti. \textbf{QED}
    }%
}

\textbf{Oss.} questo teorema implica non solo che lo stato fondamentale non ha mai nodi, ma anche che gli
stati eccitati ne abbiano: considerando il primo stato eccitato $\psi_1(x)$, deve valere che 
$0 = \bra{\psi_0}\ket{\psi_1}=\int_{}^{}dx\psi^*_0(x)\psi_1(x)$, che non è possibile se $\psi_1(x)$ non
si annulla almeno una volta; analogamente si ragiona per $\psi_2(x)$ e si trova che l'n-esimo stato
eccitato deve avere almeno $n$ nodi.

\noindent\fbox{%
    \parbox{\textwidth}{%
    \textit{Corollario} Lo stato fondamentale è autostato di tutti gli operatori che commutano con $H$.

    \textit{Dimostrazione} Dato $F/[H,F]=0$, si ha
    $$ HF\ket{\psi_0} = FH\ket{\psi_0} = E_0F\ket{\psi_0}$$
    che significa che anche lo stato $F\ket{\psi_0}$ ha autovalore $E_0$, dunque per il teorema 
    $F\ket{\psi_0} \propto\ket{\psi_0}$, cioè 
    $F\ket{\psi_0} = \lambda\ket{\psi_0}$
    \textbf{QED}
    }%
}
\newpage
\textbf{Oss.} questo è un risultato molto importante, perché asserisce che lo stato fondamentale di $H$
presenta le stesse simmetrie di $H$ che, in altre parole, significa che \textbf{in Meccanica Quantistica
non esiste la rottura spontanea di simmetria}, quando i gradi di libertà sono finiti.

Questo non è vero in meccanica classica, si prenda come esempio il cosiddetto potenziale a \textit{sombrero},
che è una funzione pari: i moti con energia minima sono banali (la particella rimane ferma o a destra
o a sinistra del picco centrale della funzione) e non presentano parità, in altre parole 
\textbf{viene rotta la simmetria di parità del potenziale}.

\begin{figure}[H]
    \centering
    \includegraphics[width = \linewidth]{./sombrero.png}
    \caption{Potenziale $V(x) = x^4 - x^2$, un esempio di potenziale a \textit{sombrero}}
\end{figure}

In meccanica quantistica, invece, lo stato fondamentale di un potenziale del genere può fare
\textit{tunneling} e, per il teorema precedente, \textbf{deve} essere pari.

Un altro esempio manifesto di questo risultato è lo stato fondamentale dell'atomo di idrogeno: l'hamiltoniana
coulombiana è invariante per rotazioni attorno a un qualunque asse, infatti lo stato fondamentale
$\psi_{100}(r,\theta,\varphi)$ ha simmetria sferica (poiché la ha l'armonica sferica
$Y_0^0(\theta,\varphi)$)


La rottura spontanea di simmetria è possibile invece in teorie quantistiche a \textit{infiniti} gradi
di libertà, cioè nelle teorie di campo quantistiche.

\newpage
\section{Prodotto diretto tra Spazi di Hilbert}
Nel corso di Meccanica Quantistica 1 si sono incontrati degli stati dipendenti da 
più numeri quantici, legati agli autovalori di un SCOC: in particolare ci sono due 
diverse situazioni
\begin{itemize}
    \item Gli autovalori del SCOC \textbf{non sono indipendenti} tra loro, come
	nel caso di $j,m$ dei momenti angolari;
    \item Gli autovalori del SCOC \textbf{sono indipendenti} tra loro, come il
	numero quantico di spin $s$ e il numero quantico dell'energia $n$.
\end{itemize}
Nel secondo caso, i gradi di libertà si possono considerare indipendenti, cioè gli
autovettori appartengono a spazi di Hilbert \textit{diversi}. 

In seguito è presentata la trattazione formale di questo concetto.
\subsection{Costruzione dello spazio prodotto}
\textbf{Def.} Dati due vettori in spazi di Hilbert
$\ket{a}_1\in\mathcal{H}_1,\ket{b}_2\in\mathcal{H}_2$, il \textit{prodotto diretto}
(o tensoriale) $\ket{a}_1\otimes\ket{b}_2 := \ket{c}_{\otimes}$ è un'applicazione
lineare e continua dell'insieme di omomorfsmi 
$\text{Hom}(\mathcal{H}_1^*,\mathcal{H}_2) \sim \text{Hom}(\mathcal{H}_2^*,\mathcal{H}_1)$.

Questo significa che il prodotto tensoriale è, contemporaneamente
\begin{itemize}
    \item $\ket{c}_\otimes: \mathcal{H}_1^*\to\mathcal{H}_2$, cioè tale che
	$_1\bra{d}\ket{c}_\otimes := ({}_1\bra{d}\ket{a}_1)\ket{b}_2\in\mathcal{H}_2$;
    \item $\ket{c}_\otimes: \mathcal{H}_2^*\to\mathcal{H}_1$, cioè tale che
	$_2\bra{e}\ket{c}_\otimes := (_2\bra{e}\ket{b}_2)\ket{a}_1\in\mathcal{H}_1$.
\end{itemize}

Bisogna inoltre richiedere che $\otimes$ sia associativo rispetto alle strutture
lineari di $\mathcal{H}_1$ e di $\mathcal{H}_2$:
$(\lambda\ket{a}_1 + \mu\ket{b}_1)\otimes\ket{c}_2 =
\lambda\ket{a}_1\otimes\ket{c}_2 + \mu\ket{b}_1\otimes\ket{c}_2$ e viceversa.

Si noti però che ${\ket{c}_\otimes}$ non è uno spazio vettoriale, poiché non è detto
che esistano due vettori $\ket{\alpha}_1,\ket{\beta}_2$ tali che
$\ket{a}_1\otimes\ket{b}_2 + \ket{c}_1\otimes\ket{d}_2 =
\ket{\alpha}_1\otimes\ket{\beta}_2$: il fine è di costruire un nuovo spazio di
Hilbert, dunque le combinazioni lineari di prodotti tensoriali dovranno essere 
aggiunte "a forza" nello spazio che si vuole creare.

Lo spazio prodotto è unitario rispetto al prodotto scalare
$_\otimes\bra{c_1}\ket{c_2}_\otimes :=
(_1\bra{a_1}\ket{a_2}_1) (_2\bra{b_1}\ket{b_2}_2)$, che gli conferisce una struttura
metrica e topologica, come già discusso.

L'ultimo requisito da soddisfare è la completezza: come per la struttura vettoriale,
si aggiungono allo spazio tutti i limiti delle successioni e delle serie di Cauchy 
in modo da renderlo completo.

\textbf{Def.} Lo spazio prodotto diretto (o tensoriale) di due spazi di Hilbert 
$\mathcal{H}_1,\mathcal{H}_2$ è l'insieme dei prodotti tensoriali dei loro elementi,
unito alle loro combinazioni lineari e ai limiti di successioni e serie di Cauchy.
Si scrive $\mathcal{H}=\mathcal{H}_1\otimes\mathcal{H}_2$ ed è uno spazio di Hilbert.
\subsection{Base dello spazio prodotto}
Il generico elemento di $\mathcal{H}$, per la completezza dello spazio, si può
scrivere come
$$\ket{v}=\sum_{k}^{\infty} \Lambda_k\ket{a_k}_1\otimes\ket{b_k}_2$$
con $\Lambda_k$ coefficienti complessi; scelte delle basi per $\mathcal{H}_1
,\mathcal{H}_2$, si può scrivere

$$\ket{a_k}_1 = \sum_{i}^{}c_{ki}\ket{i}_1 \text{ e }
\ket{b_k}_2 = \sum_{j}^{}d_{kj}\ket{j}_2$$ 
e allora si ha
$$
\ket{v} = \sum_{k,i,j}^{\infty} \Lambda_k c_{ki}d_{kj}\ket{i}_1\otimes\ket{j}_2 :=
\sum_{i,j}^{\infty} \lambda_{ij}\ket{i}_1\otimes\ket{j}_2
$$
con $\lambda_{ij}:=\sum_{k}^{\infty}\Lambda_k c_{ki} d_{kj}$

La base è dunque $\ket{i}_1\otimes\ket{j}_2$ ed è ortonormale e completa per costruzione,
infatti 
$$\{\bra{i}\otimes\bra{j}\}\{\ket{k}\otimes\ket{l}\} = \delta_{ik}\delta_{jl}$$
\subsection{Operatori sullo spazio prodotto}
Un operatore su $\mathcal{H}_1\otimes\mathcal{H}_2 = \mathcal{H}$ è definito come su un qualunque
spazio di Hilbert, ma può essere "ridotto" a un operatore su uno dei due spazi iniziali;

Denotando con $\ket{il}:=\ket{i}_1\otimes\ket{l}_2$ la base di $\mathcal{H}$, la rappresentazione
di un operatore $F$ su $\mathcal{H}$ rispetto alla base è

$$
F = F\mathbb{I} = \sum_{i,j,l,m}^{\infty} \ket{il}\bra{il}F\ket{jm}\bra{jm} :=
\sum_{i,j,l,m}^{\infty} \ket{il}F_{il,jm}\bra{jm}
$$
dove $F_{il,jm}\in\mathbb{C}$ sono i coefficienti che definiscono l'operatore rispetto alla base.

Dato l'operatore $F$ e due vettori $\ket{\alpha}_1,\ket{\beta}_1\in\mathcal{H}_1$, si definisce un
operatore su $\mathcal{H}_2$ nel seguente modo:
$$
_1\bra{\alpha}F\ket{\beta}_1 = \sum_{i,l,j,m}^{\infty} ({}_1\bra{\alpha}\ket{i}_1)\ket{l}_2 F_{il,jm}
\;\bra{m}_2\; ({}_1\bra{j}\ket{\beta}_1) := 
\sum_{l,m}^{\infty} \ket{l}_2 \hat{F}_{lm}\bra{m}_2
$$
che è un operatore su $\mathcal{2}$, dove
$\hat{F}_{lm} := \sum_{i,j}^{\infty} ({}_1\bra{\alpha}\ket{i}_1) F_{il,jm} ({}_1\bra{j}\ket{\beta}_1) $
è la rappresentazione dell'operatore rispetto alla base di $\mathcal{H}_2$

\textbf{Esempio} La traccia di un operatore su $\mathcal{H}_1$ è 
$$Tr_1(F) := \sum_{k}^{\infty} {}_1\bra{k}F\ket{k}_1 =
\sum_{l,m}^{\infty} \ket{l}_2 \left(\sum_{k}^{\infty} F_{kl,km} \right)\bra{m}_2$$
ed è un operatore su $\mathcal{H}_2$


Un altro modo di costruire operatori sullo spazio prodotto è definire il prodotto tensoriale tra
operatori:

\textbf{Def.} Dati $A_1,A_2$ operatori su $\mathcal{H}_1,\mathcal{H}_2$, il loro prodotto tensoriale
$A_1\otimes A_2 = B$ è un operatore su $\mathcal{H}_1\otimes\mathcal{H}_2$ tale che
$$
B(\ket{a}_1\otimes\ket{b}_2) := A_1\ket{a}_1\otimes A_2\ket{b}_2
$$


\textbf{Oss.} Per definire operatori che agiscono solo su uno dei due spazi, basta moltiplicarli
per l'identità: $B=\mathbb{I}\otimes{A_2} \to B(\ket{a}_1\otimes\ket{b}_2) =
\ket{a}_1\otimes A_2\ket{b}_2$
